 \emph{Interpretieren Sie Ihre Ergebnisse. Ist die Leistung akzeptabel? Gehen Sie auf die Ausgangsfragen zurück und vergleichen Sie mit anderen in der Literatur beschriebenen Methoden. Wie erklären Sie unerwartete Ergebnisse? Welche Unzulänglichkeiten oder Einschränkungen weisen Ihre Methoden und Ergebnisse auf? Wo gibt es Raum für Verbesserungen? Könnte eine Kombination mit Methoden aus der Literatur die Arbeit weiter verbessern? Nennen Sie Pläne, Bedürfnisse und Empfehlungen für die zukünftige Forschung. Dies kann auch in einem separaten Unterabschnitt ``Zukünftige Arbeiten'' erfolgen.}
 
 
 %---------------------------------------------------
 \subsection*{Bewertung}

%\label{sec:bewertung}

In die Bewertung der Abschlussarbeit gehen zum Einen die Bewertung der schriftlichen Abschlussarbeit an sich, die selbständige Arbeitsweise und die Bewertung des Ergebnisses und zum Anderen die Bewertung des Kolloquiums ein. Jeweils wird nach zehn gleichgewichteten, in Gruppen unterteilten Kriterien bewertet.



Jedes der zehn Kriterien wird mit einer Note zwischen $1$ und $5$ bewertet. („1“ sehr gut; „2“ gut; „3“ befriedigend; „4“ ausreichend; „5“ nicht ausreichend)  Erklärungen zu den einzelnen Noten werden bei Bedarf in Form eines kurzen Gutachtens ergänzt. (maximal eine DIN A4 Seite)

%\subsection{Bewertung der Abschlussarbeit}

\begin{table}[!ph]
\centering
\begin{tabular}{ |p{1cm}|p{10cm}|p{0.7cm}|  }
\hline
Lfd.Nr. & Kriterien  & Note\\  
\hline
\hline
&\textbf{I. Bearbeitungsphase}&\\
\hline
1&Einarbeitung, Verständnis und Überblick unter Berücksichtigung des Schwierigkeitsgrades&\\
\hline
2& Einfallsreichtum, Ideen, Initiative, Selbständigkeit, Entscheidungsfreudigkeit&\\
\hline
3& systematisches Vorgehen, wissenschaftliche Eigenleistung&\\
\hline
4&Organisationstalent, Zeiteinteilung; Ausdauer, Fleiß&\\
\hline
& \textbf{II. Ergebnis (Theorie, Software, Hardware)}&\\ 
\hline
5&Inwieweit wurde das mögliche Ziel unter Berücksichtigung der Anforderungen erreicht&\\
\hline
6&Verwendbarkeit der Ergebnisse: Tragfähigkeit der Theorie und Methodik und/oder Funktionstüchtigkeit der Hardware und der Software&\\
\hline
&\textbf{[III.] Schriftliche Ausarbeitung (Dokumentation)}&\\
\hline
7&Klarheit von Gliederung und Aufbau&\\
\hline
8&Vollständigkeit und wissenschaftliche Korrektheit, Bewertung der Ergebnisse&\\
\hline
9&Stil, Ausdruck&\\\hline
10&Prägnanz der Abbildungen und Diagramme, grafische Gestaltung&\\
\hline
\end{tabular}
\caption{Bewertungskriterien für die Abschlussarbeit.}
\label{tab:arbeit}
\end{table}

%\subsection{Bewertung des Kolloquiums}

\begin{table}[!ph]
\centering
\begin{tabular}{ |p{1cm}|p{10cm}|p{0.7cm}|  }
 \hline
 Lfd.Nr. & Kriterien & Note \\  
\hline\hline
& \textbf{I. Präsentation}&\\
\hline
1   &didaktischer Aufbau des Vortrags&\\
\hline
2   &Vortrag: fachlicher Inhalt, kritische Bewertung der Ergebnisse&\\
\hline
3   &Vortragsstil, Zeitdisziplin&\\
\hline
4   &Qualität der Vortagsmaterialien: Folien, Video, Demonstratoren etc., wissenschaftliche Korrektheit&\\
\hline
& \textbf{II. Diskussionsverhalten}&\\
\hline
5   &Diskussionsverhalten bezüglich des Vortrags&\\
\hline
6   &Diskussionsverhalten bezüglich der schriftlichen Ausarbeitung&\\
\hline
7   &Fragen und Diskussion: fachliche Kompetenz (Ingenieurstechnische Grundlagen)&\\
\hline
8   &Fragen und Diskussion: fachliche Kompetenz (Elektrotechnik)&\\
\hline
9   &Fragen und Diskussion: fachliche Kompetenz (themenspezifisch)&\\
\hline 
10  &Fragen und Diskussion: methodische Kompetenz&\\
\hline
\end{tabular}
\caption{Bewertungskriterien für das Kolloquium.}
\label{tab:kolloquium}
\end{table}


